% Options for packages loaded elsewhere
% Options for packages loaded elsewhere
\PassOptionsToPackage{unicode}{hyperref}
\PassOptionsToPackage{hyphens}{url}
\PassOptionsToPackage{dvipsnames,svgnames,x11names}{xcolor}
%
\documentclass[
  letterpaper,
  DIV=11,
  numbers=noendperiod]{scrartcl}
\usepackage{xcolor}
\usepackage{amsmath,amssymb}
\setcounter{secnumdepth}{-\maxdimen} % remove section numbering
\usepackage{iftex}
\ifPDFTeX
  \usepackage[T1]{fontenc}
  \usepackage[utf8]{inputenc}
  \usepackage{textcomp} % provide euro and other symbols
\else % if luatex or xetex
  \usepackage{unicode-math} % this also loads fontspec
  \defaultfontfeatures{Scale=MatchLowercase}
  \defaultfontfeatures[\rmfamily]{Ligatures=TeX,Scale=1}
\fi
\usepackage{lmodern}
\ifPDFTeX\else
  % xetex/luatex font selection
\fi
% Use upquote if available, for straight quotes in verbatim environments
\IfFileExists{upquote.sty}{\usepackage{upquote}}{}
\IfFileExists{microtype.sty}{% use microtype if available
  \usepackage[]{microtype}
  \UseMicrotypeSet[protrusion]{basicmath} % disable protrusion for tt fonts
}{}
\makeatletter
\@ifundefined{KOMAClassName}{% if non-KOMA class
  \IfFileExists{parskip.sty}{%
    \usepackage{parskip}
  }{% else
    \setlength{\parindent}{0pt}
    \setlength{\parskip}{6pt plus 2pt minus 1pt}}
}{% if KOMA class
  \KOMAoptions{parskip=half}}
\makeatother
% Make \paragraph and \subparagraph free-standing
\makeatletter
\ifx\paragraph\undefined\else
  \let\oldparagraph\paragraph
  \renewcommand{\paragraph}{
    \@ifstar
      \xxxParagraphStar
      \xxxParagraphNoStar
  }
  \newcommand{\xxxParagraphStar}[1]{\oldparagraph*{#1}\mbox{}}
  \newcommand{\xxxParagraphNoStar}[1]{\oldparagraph{#1}\mbox{}}
\fi
\ifx\subparagraph\undefined\else
  \let\oldsubparagraph\subparagraph
  \renewcommand{\subparagraph}{
    \@ifstar
      \xxxSubParagraphStar
      \xxxSubParagraphNoStar
  }
  \newcommand{\xxxSubParagraphStar}[1]{\oldsubparagraph*{#1}\mbox{}}
  \newcommand{\xxxSubParagraphNoStar}[1]{\oldsubparagraph{#1}\mbox{}}
\fi
\makeatother

\usepackage{color}
\usepackage{fancyvrb}
\newcommand{\VerbBar}{|}
\newcommand{\VERB}{\Verb[commandchars=\\\{\}]}
\DefineVerbatimEnvironment{Highlighting}{Verbatim}{commandchars=\\\{\}}
% Add ',fontsize=\small' for more characters per line
\usepackage{framed}
\definecolor{shadecolor}{RGB}{241,243,245}
\newenvironment{Shaded}{\begin{snugshade}}{\end{snugshade}}
\newcommand{\AlertTok}[1]{\textcolor[rgb]{0.68,0.00,0.00}{#1}}
\newcommand{\AnnotationTok}[1]{\textcolor[rgb]{0.37,0.37,0.37}{#1}}
\newcommand{\AttributeTok}[1]{\textcolor[rgb]{0.40,0.45,0.13}{#1}}
\newcommand{\BaseNTok}[1]{\textcolor[rgb]{0.68,0.00,0.00}{#1}}
\newcommand{\BuiltInTok}[1]{\textcolor[rgb]{0.00,0.23,0.31}{#1}}
\newcommand{\CharTok}[1]{\textcolor[rgb]{0.13,0.47,0.30}{#1}}
\newcommand{\CommentTok}[1]{\textcolor[rgb]{0.37,0.37,0.37}{#1}}
\newcommand{\CommentVarTok}[1]{\textcolor[rgb]{0.37,0.37,0.37}{\textit{#1}}}
\newcommand{\ConstantTok}[1]{\textcolor[rgb]{0.56,0.35,0.01}{#1}}
\newcommand{\ControlFlowTok}[1]{\textcolor[rgb]{0.00,0.23,0.31}{\textbf{#1}}}
\newcommand{\DataTypeTok}[1]{\textcolor[rgb]{0.68,0.00,0.00}{#1}}
\newcommand{\DecValTok}[1]{\textcolor[rgb]{0.68,0.00,0.00}{#1}}
\newcommand{\DocumentationTok}[1]{\textcolor[rgb]{0.37,0.37,0.37}{\textit{#1}}}
\newcommand{\ErrorTok}[1]{\textcolor[rgb]{0.68,0.00,0.00}{#1}}
\newcommand{\ExtensionTok}[1]{\textcolor[rgb]{0.00,0.23,0.31}{#1}}
\newcommand{\FloatTok}[1]{\textcolor[rgb]{0.68,0.00,0.00}{#1}}
\newcommand{\FunctionTok}[1]{\textcolor[rgb]{0.28,0.35,0.67}{#1}}
\newcommand{\ImportTok}[1]{\textcolor[rgb]{0.00,0.46,0.62}{#1}}
\newcommand{\InformationTok}[1]{\textcolor[rgb]{0.37,0.37,0.37}{#1}}
\newcommand{\KeywordTok}[1]{\textcolor[rgb]{0.00,0.23,0.31}{\textbf{#1}}}
\newcommand{\NormalTok}[1]{\textcolor[rgb]{0.00,0.23,0.31}{#1}}
\newcommand{\OperatorTok}[1]{\textcolor[rgb]{0.37,0.37,0.37}{#1}}
\newcommand{\OtherTok}[1]{\textcolor[rgb]{0.00,0.23,0.31}{#1}}
\newcommand{\PreprocessorTok}[1]{\textcolor[rgb]{0.68,0.00,0.00}{#1}}
\newcommand{\RegionMarkerTok}[1]{\textcolor[rgb]{0.00,0.23,0.31}{#1}}
\newcommand{\SpecialCharTok}[1]{\textcolor[rgb]{0.37,0.37,0.37}{#1}}
\newcommand{\SpecialStringTok}[1]{\textcolor[rgb]{0.13,0.47,0.30}{#1}}
\newcommand{\StringTok}[1]{\textcolor[rgb]{0.13,0.47,0.30}{#1}}
\newcommand{\VariableTok}[1]{\textcolor[rgb]{0.07,0.07,0.07}{#1}}
\newcommand{\VerbatimStringTok}[1]{\textcolor[rgb]{0.13,0.47,0.30}{#1}}
\newcommand{\WarningTok}[1]{\textcolor[rgb]{0.37,0.37,0.37}{\textit{#1}}}

\usepackage{longtable,booktabs,array}
\usepackage{calc} % for calculating minipage widths
% Correct order of tables after \paragraph or \subparagraph
\usepackage{etoolbox}
\makeatletter
\patchcmd\longtable{\par}{\if@noskipsec\mbox{}\fi\par}{}{}
\makeatother
% Allow footnotes in longtable head/foot
\IfFileExists{footnotehyper.sty}{\usepackage{footnotehyper}}{\usepackage{footnote}}
\makesavenoteenv{longtable}
\usepackage{graphicx}
\makeatletter
\newsavebox\pandoc@box
\newcommand*\pandocbounded[1]{% scales image to fit in text height/width
  \sbox\pandoc@box{#1}%
  \Gscale@div\@tempa{\textheight}{\dimexpr\ht\pandoc@box+\dp\pandoc@box\relax}%
  \Gscale@div\@tempb{\linewidth}{\wd\pandoc@box}%
  \ifdim\@tempb\p@<\@tempa\p@\let\@tempa\@tempb\fi% select the smaller of both
  \ifdim\@tempa\p@<\p@\scalebox{\@tempa}{\usebox\pandoc@box}%
  \else\usebox{\pandoc@box}%
  \fi%
}
% Set default figure placement to htbp
\def\fps@figure{htbp}
\makeatother





\setlength{\emergencystretch}{3em} % prevent overfull lines

\providecommand{\tightlist}{%
  \setlength{\itemsep}{0pt}\setlength{\parskip}{0pt}}



 


\usepackage{fvextra}
\DefineVerbatimEnvironment{Highlighting}{Verbatim}{breaklines,commandchars=\\\{\}}
\KOMAoption{captions}{tableheading}
\makeatletter
\@ifpackageloaded{caption}{}{\usepackage{caption}}
\AtBeginDocument{%
\ifdefined\contentsname
  \renewcommand*\contentsname{Table of contents}
\else
  \newcommand\contentsname{Table of contents}
\fi
\ifdefined\listfigurename
  \renewcommand*\listfigurename{List of Figures}
\else
  \newcommand\listfigurename{List of Figures}
\fi
\ifdefined\listtablename
  \renewcommand*\listtablename{List of Tables}
\else
  \newcommand\listtablename{List of Tables}
\fi
\ifdefined\figurename
  \renewcommand*\figurename{Figure}
\else
  \newcommand\figurename{Figure}
\fi
\ifdefined\tablename
  \renewcommand*\tablename{Table}
\else
  \newcommand\tablename{Table}
\fi
}
\@ifpackageloaded{float}{}{\usepackage{float}}
\floatstyle{ruled}
\@ifundefined{c@chapter}{\newfloat{codelisting}{h}{lop}}{\newfloat{codelisting}{h}{lop}[chapter]}
\floatname{codelisting}{Listing}
\newcommand*\listoflistings{\listof{codelisting}{List of Listings}}
\makeatother
\makeatletter
\makeatother
\makeatletter
\@ifpackageloaded{caption}{}{\usepackage{caption}}
\@ifpackageloaded{subcaption}{}{\usepackage{subcaption}}
\makeatother
\makeatletter
\definecolor{QuartoInternalColor1}{rgb}{0,0,0}
\makeatother
\makeatletter
\@ifpackageloaded{tikz}{}{\usepackage{tikz}}
\makeatother
        \newcommand*\circled[1]{\tikz[baseline=(char.base)]{
          \node[shape=circle,draw,inner sep=1pt] (char) {{\scriptsize#1}};}}  
                  
\usepackage{bookmark}
\IfFileExists{xurl.sty}{\usepackage{xurl}}{} % add URL line breaks if available
\urlstyle{same}
\hypersetup{
  pdftitle={Homework 3 Solutions},
  colorlinks=true,
  linkcolor={blue},
  filecolor={Maroon},
  citecolor={Blue},
  urlcolor={Blue},
  pdfcreator={LaTeX via pandoc}}


\title{Homework 3 Solutions}
\usepackage{etoolbox}
\makeatletter
\providecommand{\subtitle}[1]{% add subtitle to \maketitle
  \apptocmd{\@title}{\par {\large #1 \par}}{}{}
}
\makeatother
\subtitle{BEE 4850/5850}
\author{}
\date{}
\begin{document}
\maketitle

\RecustomVerbatimEnvironment{verbatim}{Verbatim}{
showspaces = false,
showtabs = false,
breaksymbolleft={},
breaklines
% Note: setting commandchars=\\\{\} here will cause an error
}


\subsection{Overview}\label{overview}

\subsubsection{Instructions}\label{instructions}

The goal of this homework assignment is to practice developing and
working with probability models for data.

\begin{itemize}
\tightlist
\item
  Problem 1 asks you to diagnose autocorrelation and develop an AR model
  for the weather-based variability at the Sewell's Point tide gauge.
\item
  Problem 2 asks you to fit a sea-level rise model under both
  assumptions of Gaussian i.i.d. residuals and autocorrelated residuals.
\item
  Problem 3 asks you to conduct an analysis of extreme temperatures in
  Ithaca under both stationary and non-stationary assumptions.
\item
  Problem 4 asks you to model how often Ithaca temperatures exceed
  thresholds.
\end{itemize}

\subsubsection{Load Environment}\label{load-environment}

The following code loads the environment and makes sure all needed
packages are installed. This should be at the start of most Julia
scripts.

\begin{Shaded}
\begin{Highlighting}[]
\ImportTok{import} \BuiltInTok{Pkg}
\BuiltInTok{Pkg}\NormalTok{.}\FunctionTok{activate}\NormalTok{(}\PreprocessorTok{@\_\_DIR\_\_}\NormalTok{)}
\BuiltInTok{Pkg}\NormalTok{.}\FunctionTok{instantiate}\NormalTok{()}
\end{Highlighting}
\end{Shaded}

The following packages are included in the environment (to help you find
other similar packages in other languages). The code below loads these
packages for use in the subsequent notebook (the desired functionality
for each package is commented next to the package).

\begin{Shaded}
\begin{Highlighting}[]
\ImportTok{using} \BuiltInTok{Random} \CommentTok{\# random number generation and seed{-}setting}
\ImportTok{using} \BuiltInTok{DataFrames} \CommentTok{\# tabular data structure}
\ImportTok{using} \BuiltInTok{DataFramesMeta}
\ImportTok{using} \BuiltInTok{CSV} \CommentTok{\# reads/writes .csv files}
\ImportTok{using} \BuiltInTok{Distributions} \CommentTok{\# interface to work with probability distributions}
\ImportTok{using} \BuiltInTok{Plots} \CommentTok{\# plotting library}
\ImportTok{using} \BuiltInTok{StatsBase} \CommentTok{\# statistical quantities like mean, median, etc}
\ImportTok{using} \BuiltInTok{StatsPlots} \CommentTok{\# some additional statistical plotting tools}
\ImportTok{using} \BuiltInTok{Optim} \CommentTok{\# optimization tools}
\ImportTok{using} \BuiltInTok{Dates} \CommentTok{\# date{-}time interface}

\BuiltInTok{Random}\NormalTok{.}\FunctionTok{seed!}\NormalTok{(}\FloatTok{1}\NormalTok{)}
\end{Highlighting}
\end{Shaded}

\subsection{Problems}\label{problems}

\subsubsection{Scoring}\label{scoring}

\begin{itemize}
\tightlist
\item
  Problem 1 is worth 5 points.
\item
  Problem 2 is worth 8 points.
\item
  Problem 3 is worth 7 points.
\item
  Problem 4 is worth 5 points.
\end{itemize}

\subsubsection{Problem 1}\label{problem-1}

\paragraph{Problem 1.1}\label{problem-1.1}

To load and process the data:

\begin{Shaded}
\begin{Highlighting}[]
\NormalTok{tide\_dat }\OperatorTok{=}\NormalTok{ CSV.}\FunctionTok{read}\NormalTok{(}\FunctionTok{joinpath}\NormalTok{(}\StringTok{"data"}\NormalTok{, }\StringTok{"norfolk{-}hourly{-}surge{-}2015.csv"}\NormalTok{), DataFrame)}
\NormalTok{surge\_resids }\OperatorTok{=}\NormalTok{ tide\_dat[}\OperatorTok{:}\NormalTok{, }\FloatTok{5}\NormalTok{] }\OperatorTok{{-}}\NormalTok{ tide\_dat[}\OperatorTok{:}\NormalTok{, }\FloatTok{3}\NormalTok{]}
\FunctionTok{plot}\NormalTok{(surge\_resids; ylabel}\OperatorTok{=}\StringTok{"Gauge Measurement (m)"}\NormalTok{, label}\OperatorTok{=}\StringTok{"Weather{-}Related Residual"}\NormalTok{, legend}\OperatorTok{=:}\NormalTok{topleft, xlabel}\OperatorTok{=}\StringTok{"Hour Number"}\NormalTok{)}
\end{Highlighting}
\end{Shaded}

\begin{figure}[H]

\centering{

\pandocbounded{\includegraphics[keepaspectratio]{hw03_files/mediabag/hw03_files/figure-pdf/fig-p1-data-output-1.pdf}}

}

\caption{\label{fig-p1-data}Plot of the weather-related variability from
the 2015 Sewell's Point hourly tide gauge data.}

\end{figure}%

Let's look at the partial autocorrelation function of the residuals to
determine what the autoregressive lag should be.

\begin{Shaded}
\begin{Highlighting}[]
\NormalTok{resid\_acf }\OperatorTok{=} \FunctionTok{pacf}\NormalTok{(surge\_resids, }\FloatTok{1}\OperatorTok{:}\FloatTok{5}\NormalTok{)}
\FunctionTok{plot}\NormalTok{(resid\_acf , marker}\OperatorTok{=:}\NormalTok{circle, line}\OperatorTok{=:}\NormalTok{stem, linewidth}\OperatorTok{=}\FloatTok{3}\NormalTok{, markersize}\OperatorTok{=}\FloatTok{5}\NormalTok{, legend}\OperatorTok{=}\ConstantTok{false}\NormalTok{, ylabel}\OperatorTok{=}\StringTok{"Partial Autocorrelation"}\NormalTok{, xlabel}\OperatorTok{=}\StringTok{"Time Lag"}\NormalTok{)}
\end{Highlighting}
\end{Shaded}

\begin{figure}[H]

\centering{

\pandocbounded{\includegraphics[keepaspectratio]{hw03_files/mediabag/hw03_files/figure-pdf/fig-p1-acf-output-1.pdf}}

}

\caption{\label{fig-p1-acf}Autocorrelation plot of the tide gauge
residuals.}

\end{figure}%

There is clearly a strong lag-1 autocorrelation in
Figure~\ref{fig-p1-acf}, so we would need to at least use an AR(1)
model. You could also argue that we should include a second lag, as
well, as the lag-2 partial autocorrelation is non-negligible, but this
is less critical than the AR(1).

\paragraph{Problem 1.2}\label{problem-1.2}

This AR(1) model is \[\begin{aligned}
y_t &= \alpha + \rho y_{t+1} + \varepsilon \\
\varepsilon &\sim N(0, \sigma) \\
y_0 &\sim N\left(\frac{\alpha}{1 - \rho}, \frac{\sigma^2}{1-\rho^2}\right).
\end{aligned}\]

Let's implement this model and fit it. We'll specify the likelihood
based on whitening the residuals and computing the individual normal
likelihoods of the resulting terms, but you should get the same
parameters if you used the joint likelihood of the entire time series.

\phantomsection\label{annotated-cell-5}%
\begin{Shaded}
\begin{Highlighting}[]
\KeywordTok{function} \FunctionTok{llik\_ar}\NormalTok{(params, y)}
\NormalTok{    a, p, s }\OperatorTok{=}\NormalTok{ params }\hspace*{\fill}\NormalTok{\circled{1}}
\NormalTok{    ll }\OperatorTok{=} \FloatTok{0} \CommentTok{\# initialize log{-}likelihood counter}
\NormalTok{    ll }\OperatorTok{+=} \FunctionTok{logpdf}\NormalTok{(}\FunctionTok{Normal}\NormalTok{(a }\OperatorTok{/}\NormalTok{ (}\FloatTok{1} \OperatorTok{{-}}\NormalTok{ p), s}\OperatorTok{/}\FunctionTok{sqrt}\NormalTok{(}\FloatTok{1}\OperatorTok{{-}}\NormalTok{p}\OperatorTok{\^{}}\FloatTok{2}\NormalTok{)), y[}\FloatTok{1}\NormalTok{]) }\CommentTok{\# Normal() in Distributions.jl is parameterized with the standard deviation, not the variance}
    \ControlFlowTok{for}\NormalTok{ i }\OperatorTok{=} \FloatTok{1}\OperatorTok{:}\FunctionTok{length}\NormalTok{(y)}\OperatorTok{{-}}\FloatTok{1}
\NormalTok{        residual\_whitened }\OperatorTok{=}\NormalTok{ y[i}\OperatorTok{+}\FloatTok{1}\NormalTok{] }\OperatorTok{{-}}\NormalTok{ (a }\OperatorTok{+}\NormalTok{ p }\OperatorTok{*}\NormalTok{ y[i]) }
\NormalTok{        ll }\OperatorTok{+=} \FunctionTok{logpdf}\NormalTok{(}\FunctionTok{Normal}\NormalTok{(}\FloatTok{0}\NormalTok{, s), residual\_whitened)}
    \ControlFlowTok{end}
    \ControlFlowTok{return}\NormalTok{ ll}
\KeywordTok{end}

\CommentTok{\# set up lower and upper bounds for the parameters for the optimization}
\NormalTok{lbds }\OperatorTok{=}\NormalTok{ [}\FloatTok{0.0}\NormalTok{, }\FloatTok{0.0}\NormalTok{, }\FloatTok{0.1}\NormalTok{]}
\NormalTok{ubds }\OperatorTok{=}\NormalTok{ [}\FloatTok{5.0}\NormalTok{, }\FloatTok{1.0}\NormalTok{, }\FloatTok{5.0}\NormalTok{]}
\NormalTok{p0 }\OperatorTok{=}\NormalTok{ [}\FloatTok{1.0}\NormalTok{, }\FloatTok{0.4}\NormalTok{, }\FloatTok{0.5}\NormalTok{]}
\NormalTok{optim\_out }\OperatorTok{=}\NormalTok{ Optim.}\FunctionTok{optimize}\NormalTok{(p }\OperatorTok{{-}\textgreater{}} \FunctionTok{{-}llik\_ar}\NormalTok{(p, surge\_resids), lbds, ubds, p0)}
\NormalTok{p\_tide\_mle }\OperatorTok{=}\NormalTok{ optim\_out.minimizer}
\PreprocessorTok{@show} \FunctionTok{round}\NormalTok{(p\_tide\_mle; digits}\OperatorTok{=}\FloatTok{2}\NormalTok{)}
\end{Highlighting}
\end{Shaded}

\begin{description}
\tightlist
\item[\circled{1}]
This syntax is useful as it lets us unpack a vector or tuple into named
variables, which increases readability of the subsequent code (versus
only referring to e.g.~\texttt{p{[}1{]}}).
\end{description}

\begin{verbatim}
┌ Warning: Terminated early due to NaN in gradient.
└ @ Optim ~/.julia/packages/Optim/DtV5C/src/multivariate/optimize/optimize.jl:136
\end{verbatim}

\begin{Highlighting}
\textcolor{black}{MethodError: MethodError(Core.kwcall, ((digits = 2,), round, [0.000862808500856635, 0.9920243071864885, 0.1], RoundingMode\{:Nearest\}()), 0x0000000000006940)}
\textcolor{black}{MethodError: no method matching round(::Vector\{Float64\}, ::RoundingMode\{:Nearest\}; digits::Int64)}
\textcolor{black}{The function `round` exists, but no method is defined for this combination of argument types.}
\textcolor{black}{}\textcolor{QuartoInternalColor1}{Closest candidates are:}
\textcolor{QuartoInternalColor1}{}\textcolor{QuartoInternalColor1}{  round(}\textcolor{QuartoInternalColor1}{::Type\{T\}}\textcolor{QuartoInternalColor1}{, ::Any, }\textcolor{QuartoInternalColor1}{::RoundingMode}\textcolor{QuartoInternalColor1}{) where T>:Missing}\textcolor{QuartoInternalColor1}{ got unsupported keyword argument "digits"}\textcolor{QuartoInternalColor1}{}
\textcolor{QuartoInternalColor1}{}\textcolor{QuartoInternalColor1}{}\textcolor{QuartoInternalColor1}{   @}\textcolor{QuartoInternalColor1}{ }\textcolor{QuartoInternalColor1}{Base}\textcolor{QuartoInternalColor1}{ }\textcolor{QuartoInternalColor1}{}\textcolor{QuartoInternalColor1}{missing.jl:148}\textcolor{QuartoInternalColor1}{}\textcolor{QuartoInternalColor1}{}
\textcolor{QuartoInternalColor1}{}\textcolor{QuartoInternalColor1}{  round(}\textcolor{QuartoInternalColor1}{::Type\{T\}}\textcolor{QuartoInternalColor1}{, ::Any, }\textcolor{QuartoInternalColor1}{::RoundingMode}\textcolor{QuartoInternalColor1}{) where T}\textcolor{QuartoInternalColor1}{ got unsupported keyword argument "digits"}\textcolor{QuartoInternalColor1}{}
\textcolor{QuartoInternalColor1}{}\textcolor{QuartoInternalColor1}{}\textcolor{QuartoInternalColor1}{   @}\textcolor{QuartoInternalColor1}{ }\textcolor{QuartoInternalColor1}{Base}\textcolor{QuartoInternalColor1}{ }\textcolor{QuartoInternalColor1}{}\textcolor{QuartoInternalColor1}{rounding.jl:479}\textcolor{QuartoInternalColor1}{}\textcolor{QuartoInternalColor1}{}
\textcolor{QuartoInternalColor1}{}\textcolor{QuartoInternalColor1}{  round(}\textcolor{QuartoInternalColor1}{::Missing}\textcolor{QuartoInternalColor1}{, ::RoundingMode; sigdigits, digits, base)}
\textcolor{QuartoInternalColor1}{}\textcolor{QuartoInternalColor1}{}\textcolor{QuartoInternalColor1}{   @}\textcolor{QuartoInternalColor1}{ }\textcolor{QuartoInternalColor1}{Base}\textcolor{QuartoInternalColor1}{ }\textcolor{QuartoInternalColor1}{}\textcolor{QuartoInternalColor1}{missing.jl:144}\textcolor{QuartoInternalColor1}{}\textcolor{QuartoInternalColor1}{}
\textcolor{QuartoInternalColor1}{}\textcolor{QuartoInternalColor1}{  ...}
\textcolor{QuartoInternalColor1}{Stacktrace:}
\textcolor{QuartoInternalColor1}{ [1] }\textcolor{QuartoInternalColor1}{}\textcolor{QuartoInternalColor1}{round}\textcolor{QuartoInternalColor1}{}\textcolor{QuartoInternalColor1}{}\textcolor{QuartoInternalColor1}{(}\textcolor{QuartoInternalColor1}{}\textcolor{QuartoInternalColor1}{x}\textcolor{QuartoInternalColor1}{::}\textcolor{QuartoInternalColor1}{Vector}\textcolor{QuartoInternalColor1}{\{Float64\}}\textcolor{QuartoInternalColor1}{; }\textcolor{QuartoInternalColor1}{kws}\textcolor{QuartoInternalColor1}{::}\textcolor{QuartoInternalColor1}{@Kwargs}\textcolor{QuartoInternalColor1}{\{digits::Int64\}}\textcolor{QuartoInternalColor1}{}\textcolor{QuartoInternalColor1}{}\textcolor{QuartoInternalColor1}{)}\textcolor{QuartoInternalColor1}{}
\textcolor{QuartoInternalColor1}{}\textcolor{QuartoInternalColor1}{   @}\textcolor{QuartoInternalColor1}{ }\textcolor{QuartoInternalColor1}{Base}\textcolor{QuartoInternalColor1}{ }\textcolor{QuartoInternalColor1}{./}\textcolor{QuartoInternalColor1}{}\textcolor{QuartoInternalColor1}{}\textcolor{QuartoInternalColor1}{rounding.jl:472}\textcolor{QuartoInternalColor1}{}\textcolor{QuartoInternalColor1}{}
\textcolor{QuartoInternalColor1}{ [2] }\textcolor{QuartoInternalColor1}{}\textcolor{QuartoInternalColor1}{macro expansion}\textcolor{QuartoInternalColor1}{}
\textcolor{QuartoInternalColor1}{}\textcolor{QuartoInternalColor1}{   @}\textcolor{QuartoInternalColor1}{ }\textcolor{QuartoInternalColor1}{}\textcolor{QuartoInternalColor1}{show.jl:1232}\textcolor{QuartoInternalColor1}{}\textcolor{QuartoInternalColor1}{}\textcolor{QuartoInternalColor1}{ [inlined]}\textcolor{QuartoInternalColor1}{}
\textcolor{QuartoInternalColor1}{ [3] top-level scope}
\textcolor{QuartoInternalColor1}{}\textcolor{QuartoInternalColor1}{   @}\textcolor{QuartoInternalColor1}{ }\textcolor{QuartoInternalColor1}{\textasciitilde{}/Teaching/BEE4850/spring2026/solutions/hw03/}\textcolor{QuartoInternalColor1}{}\textcolor{QuartoInternalColor1}{}\textcolor{QuartoInternalColor1}{hw03.qmd:160}\textcolor{QuartoInternalColor1}{}\textcolor{QuartoInternalColor1}{}
\end{Highlighting}

Now we can simulate these residuals and add them back to the NOAA model
predictions using the \texttt{mapslices()} function. The resulting
histograms of the maximum and median values is shown in
Figure~\ref{fig-p2-histogram}.

\begin{Shaded}
\begin{Highlighting}[]
\CommentTok{\# T is the length of the simulated series, N the number of simulations}
\KeywordTok{function} \FunctionTok{simulate\_ar1}\NormalTok{(a, p, s, T, N)}
\NormalTok{    y }\OperatorTok{=} \FunctionTok{zeros}\NormalTok{(T, N) }\CommentTok{\# initialize storage}
\NormalTok{    y[}\FloatTok{1}\NormalTok{, }\OperatorTok{:}\NormalTok{] }\OperatorTok{=} \FunctionTok{rand}\NormalTok{(}\FunctionTok{Normal}\NormalTok{(a }\OperatorTok{/}\NormalTok{ (}\FloatTok{1} \OperatorTok{{-}}\NormalTok{ p), s}\OperatorTok{/}\FunctionTok{sqrt}\NormalTok{(}\FloatTok{1}\OperatorTok{{-}}\NormalTok{p}\OperatorTok{\^{}}\FloatTok{2}\NormalTok{)), N) }\CommentTok{\# sample initial value}
    \CommentTok{\# simulate the remaining residuals}
    \ControlFlowTok{for}\NormalTok{ t }\OperatorTok{=} \FloatTok{2}\OperatorTok{:}\NormalTok{T}
\NormalTok{        y[t, }\OperatorTok{:}\NormalTok{] }\OperatorTok{=}\NormalTok{ a }\OperatorTok{.+}\NormalTok{ p }\OperatorTok{*}\NormalTok{ y[t}\OperatorTok{{-}}\FloatTok{1}\NormalTok{, }\OperatorTok{:}\NormalTok{] }\OperatorTok{+} \FunctionTok{rand}\NormalTok{(}\FunctionTok{Normal}\NormalTok{(}\FloatTok{0}\NormalTok{, s), N)}
    \ControlFlowTok{end}
    \ControlFlowTok{return}\NormalTok{ y}
\KeywordTok{end}

\CommentTok{\# simulate new observations}
\NormalTok{simulated\_resids }\OperatorTok{=} \FunctionTok{simulate\_ar1}\NormalTok{(p\_tide\_mle[}\FloatTok{1}\NormalTok{], p\_tide\_mle[}\FloatTok{2}\NormalTok{], p\_tide\_mle[}\FloatTok{3}\NormalTok{], }\FunctionTok{length}\NormalTok{(surge\_resids), }\FloatTok{10\_000}\NormalTok{)}
\NormalTok{simulated\_obs }\OperatorTok{=} \FunctionTok{mapslices}\NormalTok{(col }\OperatorTok{{-}\textgreater{}}\NormalTok{ col }\OperatorTok{+}\NormalTok{ tide\_dat[}\OperatorTok{:}\NormalTok{, }\FloatTok{3}\NormalTok{], simulated\_resids, dims}\OperatorTok{=}\FloatTok{1}\NormalTok{)}
\CommentTok{\# plot histogram}
\NormalTok{max\_obs }\OperatorTok{=}\NormalTok{ [}\FunctionTok{maximum}\NormalTok{(simulated\_obs[}\OperatorTok{:}\NormalTok{, n]) for n }\KeywordTok{in} \FloatTok{1}\OperatorTok{:}\FunctionTok{size}\NormalTok{(simulated\_obs, }\FloatTok{2}\NormalTok{)]}
\NormalTok{med\_obs }\OperatorTok{=}\NormalTok{ [}\FunctionTok{median}\NormalTok{(simulated\_obs[}\OperatorTok{:}\NormalTok{, n]) for n }\KeywordTok{in} \FloatTok{1}\OperatorTok{:}\FunctionTok{size}\NormalTok{(simulated\_obs, }\FloatTok{2}\NormalTok{)]}
\NormalTok{p1 }\OperatorTok{=} \FunctionTok{histogram}\NormalTok{(max\_obs, xlabel}\OperatorTok{=}\StringTok{"Maximum Simulated Gauge Level (m)"}\NormalTok{, ylabel}\OperatorTok{=}\StringTok{"Count"}\NormalTok{, label}\OperatorTok{=}\StringTok{"Simulations"}\NormalTok{)}
\FunctionTok{vline!}\NormalTok{([}\FunctionTok{maximum}\NormalTok{(tide\_dat[}\OperatorTok{:}\NormalTok{, }\FloatTok{5}\NormalTok{])], color}\OperatorTok{=:}\NormalTok{red, linestyle}\OperatorTok{=:}\NormalTok{dash, label}\OperatorTok{=}\StringTok{"Observation"}\NormalTok{)}
\NormalTok{p2 }\OperatorTok{=} \FunctionTok{histogram}\NormalTok{(med\_obs, xlabel}\OperatorTok{=}\StringTok{"Median Simulated Gauge Level (m)"}\NormalTok{, ylabel}\OperatorTok{=}\StringTok{"Count"}\NormalTok{, label}\OperatorTok{=}\StringTok{"Simulations"}\NormalTok{)}
\FunctionTok{vline!}\NormalTok{([}\FunctionTok{median}\NormalTok{(tide\_dat[}\OperatorTok{:}\NormalTok{, }\FloatTok{5}\NormalTok{])], color}\OperatorTok{=:}\NormalTok{red, linestyle}\OperatorTok{=:}\NormalTok{dash, label}\OperatorTok{=}\StringTok{"Observation"}\NormalTok{)}
\FunctionTok{plot}\NormalTok{(p1, p2, layout}\OperatorTok{=}\NormalTok{(}\FloatTok{2}\NormalTok{, }\FloatTok{1}\NormalTok{))}
\end{Highlighting}
\end{Shaded}

\begin{figure}[H]

\centering{

\pandocbounded{\includegraphics[keepaspectratio]{hw03_files/mediabag/hw03_files/figure-pdf/fig-p2-histogram-output-1.pdf}}

}

\caption{\label{fig-p2-histogram}Histogram of the maximum and median
simulated tide gauge measurements.}

\end{figure}%

The median distribution fits the data well, but the observed maximum
1.98 is below all of the simulations, which could be due to a model
structural problem. For example, an AR(2) might have been a better fit.

\subsubsection{Problem 2}\label{problem-2}

\paragraph{Problem 2.1}\label{problem-2.1}

First, load the data and normalize. We'll plot the resulting series
(though this wasn't required).

\begin{Shaded}
\begin{Highlighting}[]
\NormalTok{norm\_yrs }\OperatorTok{=} \FloatTok{1880}\OperatorTok{:}\FloatTok{1900}

\NormalTok{sl\_dat }\OperatorTok{=} \FunctionTok{DataFrame}\NormalTok{(CSV.}\FunctionTok{File}\NormalTok{(}\FunctionTok{joinpath}\NormalTok{(}\StringTok{"data"}\NormalTok{,  }\StringTok{"CSIRO\_Recons\_gmsl\_yr\_2015.csv"}\NormalTok{)))}

\FunctionTok{rename!}\NormalTok{(sl\_dat, [}\OperatorTok{:}\NormalTok{Year, }\OperatorTok{:}\NormalTok{GMSLR, }\OperatorTok{:}\NormalTok{SD]) }\CommentTok{\# rename to make columns easier to work with}
\NormalTok{sl\_dat[!, }\OperatorTok{:}\NormalTok{Year] }\OperatorTok{.{-}=} \FloatTok{0.5} \CommentTok{\# shift year to line up with years instead of being half{-}year }
\NormalTok{sl\_dat[!, }\OperatorTok{:}\NormalTok{GMSLR] }\OperatorTok{.{-}=} \FunctionTok{mean}\NormalTok{(}\FunctionTok{filter}\NormalTok{(row }\OperatorTok{{-}\textgreater{}}\NormalTok{ row.Year }\OperatorTok{∈}\NormalTok{ norm\_yrs, sl\_dat)[!, }\OperatorTok{:}\NormalTok{GMSLR]) }\CommentTok{\# rescale to be relative to 1880{-}1900 mean for consistency with temperature anomaly}

\CommentTok{\# load temperature data}
\NormalTok{temp\_dat }\OperatorTok{=} \FunctionTok{DataFrame}\NormalTok{(CSV.}\FunctionTok{File}\NormalTok{(}\FunctionTok{joinpath}\NormalTok{(}\StringTok{"data"}\NormalTok{, }\StringTok{"HadCRUT.5.0.2.0.analysis.summary\_series.global.annual.csv"}\NormalTok{)))}
\FunctionTok{rename!}\NormalTok{(temp\_dat, [}\OperatorTok{:}\NormalTok{Year, }\OperatorTok{:}\NormalTok{Temp, }\OperatorTok{:}\NormalTok{Lower, }\OperatorTok{:}\NormalTok{Upper]) }\CommentTok{\# rename to make columns easier to work with}
\FunctionTok{filter!}\NormalTok{(row }\OperatorTok{{-}\textgreater{}}\NormalTok{ row.Year }\OperatorTok{∈}\NormalTok{ sl\_dat[!, }\OperatorTok{:}\NormalTok{Year], temp\_dat) }\CommentTok{\# reduce to the same years that we have SL data for}
\NormalTok{temp\_normalize }\OperatorTok{=} \FunctionTok{mean}\NormalTok{(}\FunctionTok{filter}\NormalTok{(row }\OperatorTok{{-}\textgreater{}}\NormalTok{ row.Year }\OperatorTok{∈}\NormalTok{ norm\_yrs, temp\_dat)[!, }\OperatorTok{:}\NormalTok{Temp]) }\CommentTok{\# get renormalization to rescale temperature to 1880{-}1900 mean}
\NormalTok{temp\_dat[!, }\OperatorTok{:}\NormalTok{Temp] }\OperatorTok{.{-}=}\NormalTok{ temp\_normalize}


\NormalTok{sl\_plot }\OperatorTok{=} \FunctionTok{scatter}\NormalTok{(sl\_dat[!, }\OperatorTok{:}\NormalTok{Year], sl\_dat[!, }\OperatorTok{:}\NormalTok{GMSLR], color}\OperatorTok{=:}\NormalTok{black, label}\OperatorTok{=}\StringTok{"Observations"}\NormalTok{, ylabel}\OperatorTok{=}\StringTok{"(mm)"}\NormalTok{, xlabel}\OperatorTok{=}\StringTok{"Year"}\NormalTok{, title}\OperatorTok{=}\StringTok{"Sea Level Anomaly"}\NormalTok{)}

\NormalTok{temp\_plot }\OperatorTok{=} \FunctionTok{scatter}\NormalTok{(temp\_dat[!, }\OperatorTok{:}\NormalTok{Year], temp\_dat[!, }\OperatorTok{:}\NormalTok{Temp], color}\OperatorTok{=:}\NormalTok{black, label}\OperatorTok{=}\StringTok{"Observations"}\NormalTok{, ylabel}\OperatorTok{=}\StringTok{"(°C)"}\NormalTok{, xlabel}\OperatorTok{=}\StringTok{"Year"}\NormalTok{, title}\OperatorTok{=}\StringTok{"Temperature"}\NormalTok{)}

\FunctionTok{plot}\NormalTok{(sl\_plot, temp\_plot, layout}\OperatorTok{=}\NormalTok{(}\FloatTok{2}\NormalTok{, }\FloatTok{1}\NormalTok{))}
\end{Highlighting}
\end{Shaded}

\begin{figure}[H]

\centering{

\pandocbounded{\includegraphics[keepaspectratio]{hw03_files/mediabag/hw03_files/figure-pdf/fig-p2-data-output-1.pdf}}

}

\caption{\label{fig-p2-data}Global mean sea level and temperature data
for Problem 2.}

\end{figure}%

\paragraph{Problem 2.2}\label{problem-2.2}

The discretized model is:

\[
\begin{aligned}
S(t+1) &= S(t) + \Delta t \frac{S_\text{eq}(t) - S(t)}{\tau} \\
S\text{eq}(t) &= a T(t) + b,
\end{aligned}
\]

where we take \(\delta t = 1\). Then in Julia (as an example), this
might look like the following function.

\phantomsection\label{annotated-cell-8}%
\begin{Shaded}
\begin{Highlighting}[]
\KeywordTok{function} \FunctionTok{grinsted\_slr}\NormalTok{(params, temps; Δt}\OperatorTok{=}\FloatTok{1}\NormalTok{)}
\NormalTok{    a, b, l, S₀ }\OperatorTok{=}\NormalTok{ params }\hspace*{\fill}\NormalTok{\circled{1}}
\NormalTok{    S }\OperatorTok{=} \FunctionTok{zeros}\NormalTok{(}\FunctionTok{length}\NormalTok{(temps)) }\CommentTok{\# initialize storage}
\NormalTok{    Seq }\OperatorTok{=}\NormalTok{ a }\OperatorTok{*}\NormalTok{ temps }\OperatorTok{.+}\NormalTok{ b}
\NormalTok{    S[}\FloatTok{1}\NormalTok{] }\OperatorTok{=}\NormalTok{ S₀}
    \ControlFlowTok{for}\NormalTok{ i }\OperatorTok{=} \FloatTok{2}\OperatorTok{:}\FunctionTok{length}\NormalTok{(S)}
\NormalTok{        S[i] }\OperatorTok{=}\NormalTok{ S[i}\OperatorTok{{-}}\FloatTok{1}\NormalTok{] }\OperatorTok{+}\NormalTok{  Δt }\OperatorTok{*}\NormalTok{ (Seq[i] }\OperatorTok{{-}}\NormalTok{ S[i}\OperatorTok{{-}}\FloatTok{1}\NormalTok{]) }\OperatorTok{/}\NormalTok{ l}
    \ControlFlowTok{end}
    \ControlFlowTok{return}\NormalTok{ S[}\FloatTok{1}\OperatorTok{:}\KeywordTok{end}\NormalTok{]}
\KeywordTok{end}
\end{Highlighting}
\end{Shaded}

\begin{description}
\tightlist
\item[\circled{1}]
This syntax is useful as it lets us unpack a vector or tuple into named
variables, which increases readability of the subsequent code (versus
only referring to e.g.~\texttt{p{[}1{]}}).
\end{description}

\begin{verbatim}
grinsted_slr (generic function with 1 method)
\end{verbatim}

The next function lets us calculate the log-likelihood if we assume
i.i.d. Gaussian residuals. Note the inclusion of the observation errors

\phantomsection\label{annotated-cell-9}%
\begin{Shaded}
\begin{Highlighting}[]
\KeywordTok{function} \FunctionTok{gaussian\_iid\_loglik}\NormalTok{(params, temp\_dat, slr\_obs, Δt}\OperatorTok{=}\FloatTok{1.0}\NormalTok{)}
\NormalTok{    a, b, l, S₀, s }\OperatorTok{=}\NormalTok{ params }
\NormalTok{    slr\_sim }\OperatorTok{=} \FunctionTok{grinsted\_slr}\NormalTok{((a, b, l, S₀), temp\_dat; Δt }\OperatorTok{=}\NormalTok{ Δt)}
\NormalTok{    ll }\OperatorTok{=} \FunctionTok{sum}\NormalTok{(}\FunctionTok{logpdf}\NormalTok{.(}\FunctionTok{Normal}\NormalTok{.(slr\_sim, s), slr\_obs)) }\hspace*{\fill}\NormalTok{\circled{1}}
    \ControlFlowTok{return}\NormalTok{ ll}
\KeywordTok{end}
\end{Highlighting}
\end{Shaded}

\begin{description}
\tightlist
\item[\circled{1}]
This line makes \emph{heavy} use of broadcasting. A loop would work just
as well.
\end{description}

\begin{verbatim}
gaussian_iid_loglik (generic function with 2 methods)
\end{verbatim}

Now, let's find the MLE by optimizing the \texttt{gaussian\_iid\_loglik}
function.

\begin{Shaded}
\begin{Highlighting}[]
\NormalTok{low\_bds }\OperatorTok{=}\NormalTok{ [}\FloatTok{100.0}\NormalTok{, }\FloatTok{100.0}\NormalTok{, }\FloatTok{0.1}\NormalTok{, }\OperatorTok{{-}}\FloatTok{100.0}\NormalTok{, }\FloatTok{0.0}\NormalTok{]}
\NormalTok{up\_bds }\OperatorTok{=}\NormalTok{ [}\FloatTok{2000.0}\NormalTok{, }\FloatTok{500.0}\NormalTok{, }\FloatTok{2500.0}\NormalTok{, }\FloatTok{100.0}\NormalTok{, }\FloatTok{100.0}\NormalTok{]}
\NormalTok{p₀ }\OperatorTok{=}\NormalTok{ [}\FloatTok{1000.0}\NormalTok{, }\FloatTok{200.0}\NormalTok{, }\FloatTok{1500.0}\NormalTok{, }\FloatTok{0.0}\NormalTok{, }\FloatTok{10.0}\NormalTok{]}

\NormalTok{mle\_optim }\OperatorTok{=} \FunctionTok{optimize}\NormalTok{(p }\OperatorTok{{-}\textgreater{}} \FunctionTok{{-}gaussian\_iid\_loglik}\NormalTok{(p, temp\_dat.Temp, sl\_dat.GMSLR), low\_bds, up\_bds, p₀)}
\NormalTok{p\_mle\_iid }\OperatorTok{=}\NormalTok{ mle\_optim.minimizer}
\PreprocessorTok{@show} \FunctionTok{round}\NormalTok{.(p\_mle\_iid; digits}\OperatorTok{=}\FloatTok{2}\NormalTok{)}
\end{Highlighting}
\end{Shaded}

\begin{verbatim}
round.(p_mle_iid; digits = 2) = [566.92, 222.66, 189.11, -10.36, 5.81]
\end{verbatim}

\begin{verbatim}
5-element Vector{Float64}:
 566.92
 222.66
 189.11
 -10.36
   5.81
\end{verbatim}

The sensitivity of GMSLR to increases in GMT is \(a / \tau\):

\[S(t+1) = S(t) + \frac{\overbrace{(aT + b)}^{S_\text{eq}} - S(t)}{\tau}.\]

This value is 2.5 mm/(yr\(\cdot ^\circ\)C).

Finally, the i.i.d. Gaussian model had two key assumptions:
Gaussian-distributed residuals and independent residuals. To check the
first, we can use a Gaussian Q-Q plot, and the second, a partial
autocorrelation plot.

\begin{Shaded}
\begin{Highlighting}[]
\NormalTok{grinsted\_proj\_mle }\OperatorTok{=} \FunctionTok{grinsted\_slr}\NormalTok{(p\_mle\_iid, temp\_dat.Temp)}
\NormalTok{resids }\OperatorTok{=}\NormalTok{ sl\_dat.GMSLR }\OperatorTok{{-}}\NormalTok{ grinsted\_proj\_mle}
\NormalTok{p1 }\OperatorTok{=} \FunctionTok{qqnorm}\NormalTok{(resids, ylabel}\OperatorTok{=}\StringTok{"Theoretical Values"}\NormalTok{, xlabel}\OperatorTok{=}\StringTok{"Empirical Values"}\NormalTok{)}
\NormalTok{pacf\_resids }\OperatorTok{=} \FunctionTok{pacf}\NormalTok{(resids, }\FloatTok{1}\OperatorTok{:}\FloatTok{5}\NormalTok{)}
\NormalTok{p2 }\OperatorTok{=} \FunctionTok{plot}\NormalTok{(}\FloatTok{1}\OperatorTok{:}\FloatTok{5}\NormalTok{, pacf\_resids, marker}\OperatorTok{=:}\NormalTok{circle, line}\OperatorTok{=:}\NormalTok{stem, markersize}\OperatorTok{=}\FloatTok{5}\NormalTok{, legend}\OperatorTok{=}\ConstantTok{false}\NormalTok{, ylabel}\OperatorTok{=}\StringTok{"Partial Autocorrelation"}\NormalTok{, xlabel}\OperatorTok{=}\StringTok{"Time Lag"}\NormalTok{)}

\FunctionTok{display}\NormalTok{(p1)}
\FunctionTok{display}\NormalTok{(p2)}
\end{Highlighting}
\end{Shaded}

\begin{figure}

\begin{minipage}{0.50\linewidth}

\centering{

\pandocbounded{\includegraphics[keepaspectratio]{hw03_files/mediabag/hw03_files/figure-pdf/fig-p1-residuals-output-1.pdf}}

}

\subcaption{\label{fig-p1-residuals-1}Q-Q Plot vs.~a Theoretical Normal
Distribution}

\end{minipage}%
%
\begin{minipage}{0.50\linewidth}

\centering{

\pandocbounded{\includegraphics[keepaspectratio]{hw03_files/mediabag/hw03_files/figure-pdf/fig-p1-residuals-output-2.pdf}}

}

\subcaption{\label{fig-p1-residuals-2}Partial Autocorrelation Plot}

\end{minipage}%

\caption{\label{fig-p1-residuals}Diagnostics for the model residuals.}

\end{figure}%

We can see that while the residuals appear pretty-well described by a
Gaussian distribution based on the Q-Q plot, they have a reasonable
lag-1 autocorrelation of about 0.52, so this may be an inappropriate
assumption.

\paragraph{Problem 2.3}\label{problem-2.3}

This will look like the previous problem, only with an AR(1) likelihood
for the residuals instead of the Gaussian residuals. We will assume that
the residual model is mean-zero, so \(\alpha = 0\).

\begin{Shaded}
\begin{Highlighting}[]
\KeywordTok{function} \FunctionTok{gaussian\_ar\_loglik}\NormalTok{(params, temp\_dat, slr\_obs, Δt}\OperatorTok{=}\FloatTok{1.0}\NormalTok{)}
\NormalTok{    a, b, l, S₀, p, s }\OperatorTok{=}\NormalTok{ params }
\NormalTok{    slr\_sim }\OperatorTok{=} \FunctionTok{grinsted\_slr}\NormalTok{((a, b, l, S₀), temp\_dat; Δt }\OperatorTok{=}\NormalTok{ Δt)}
\NormalTok{    slr\_resids }\OperatorTok{=}\NormalTok{ slr\_obs }\OperatorTok{{-}}\NormalTok{ slr\_sim}
\NormalTok{    ll }\OperatorTok{=} \FunctionTok{logpdf}\NormalTok{(}\FunctionTok{Normal}\NormalTok{(}\FloatTok{0}\NormalTok{, s }\OperatorTok{/} \FunctionTok{sqrt}\NormalTok{(}\FloatTok{1} \OperatorTok{{-}}\NormalTok{ p}\OperatorTok{\^{}}\FloatTok{2}\NormalTok{)), slr\_resids[}\FloatTok{1}\NormalTok{])}
    \ControlFlowTok{for}\NormalTok{ i }\OperatorTok{=} \FloatTok{2}\OperatorTok{:}\FunctionTok{length}\NormalTok{(slr\_sim)}
\NormalTok{        ll }\OperatorTok{+=} \FunctionTok{logpdf}\NormalTok{(}\FunctionTok{Normal}\NormalTok{(}\FloatTok{0}\NormalTok{, s), slr\_resids[i] }\OperatorTok{{-}}\NormalTok{ p }\OperatorTok{*}\NormalTok{ slr\_resids[i}\OperatorTok{{-}}\FloatTok{1}\NormalTok{])}
    \ControlFlowTok{end}
    \ControlFlowTok{return}\NormalTok{ ll}
\KeywordTok{end}
\end{Highlighting}
\end{Shaded}

\begin{verbatim}
gaussian_ar_loglik (generic function with 2 methods)
\end{verbatim}

Now, let's find the MLE by optimizing the \texttt{gaussian\_iid\_loglik}
function.

\begin{Shaded}
\begin{Highlighting}[]
\NormalTok{low\_bds }\OperatorTok{=}\NormalTok{ [}\FloatTok{100.0}\NormalTok{, }\FloatTok{100.0}\NormalTok{, }\FloatTok{0.1}\NormalTok{, }\OperatorTok{{-}}\FloatTok{100.0}\NormalTok{, }\FloatTok{0.0}\NormalTok{, }\FloatTok{0.1}\NormalTok{]}
\NormalTok{up\_bds }\OperatorTok{=}\NormalTok{ [}\FloatTok{2000.0}\NormalTok{, }\FloatTok{500.0}\NormalTok{, }\FloatTok{2500.0}\NormalTok{, }\FloatTok{100.0}\NormalTok{, }\FloatTok{1.0}\NormalTok{, }\FloatTok{100.0}\NormalTok{]}
\NormalTok{p₀ }\OperatorTok{=}\NormalTok{ [}\FloatTok{1000.0}\NormalTok{, }\FloatTok{200.0}\NormalTok{, }\FloatTok{1500.0}\NormalTok{, }\FloatTok{0.0}\NormalTok{, }\FloatTok{0.5}\NormalTok{, }\FloatTok{10.0}\NormalTok{]}

\NormalTok{mle\_optim }\OperatorTok{=} \FunctionTok{optimize}\NormalTok{(p }\OperatorTok{{-}\textgreater{}} \FunctionTok{{-}gaussian\_ar\_loglik}\NormalTok{(p, temp\_dat.Temp, sl\_dat.GMSLR), low\_bds, up\_bds, p₀)}
\NormalTok{p\_mle\_ar }\OperatorTok{=}\NormalTok{ mle\_optim.minimizer}
\PreprocessorTok{@show} \FunctionTok{round}\NormalTok{.(p\_mle\_ar; digits}\OperatorTok{=}\FloatTok{2}\NormalTok{)}
\end{Highlighting}
\end{Shaded}

\begin{verbatim}
round.(p_mle_ar; digits = 2) = [557.57, 216.77, 183.68, -10.44, 0.52, 4.95]
\end{verbatim}

\begin{verbatim}
6-element Vector{Float64}:
 557.57
 216.77
 183.68
 -10.44
   0.52
   4.95
\end{verbatim}

Comparing these parameters to those from Problem 2.2, all of the
parameters are a little smaller in the AR(1) residual model. The
resulting sensitivity of GMSLR to increases in GMT:

\[S(t+1) = S(t) + \frac{\overbrace{(aT + b)}^{S_\text{eq}} - S(t)}{\tau}.\]

This value is 2.6 mm/(yr\(\cdot ^\circ\)C), which is just a little
larger than with the Gaussian residuals and is unlikely to make a
meaningful difference to SLR projections; with even 4\(\cdot ^\circ\)C
of warming this would result only in an additional projected 0.4mm of
GMSLR. This can matter more in other examples, however!

\subsubsection{Problem 3}\label{problem-3}

\paragraph{Problem 3.1}\label{problem-3.1}

Let's load the data from `data/gfdl-esm4-tempmax-ithaca.csv'.

\begin{Shaded}
\begin{Highlighting}[]
\NormalTok{dat }\OperatorTok{=}\NormalTok{ CSV.}\FunctionTok{read}\NormalTok{(}\FunctionTok{joinpath}\NormalTok{(}\StringTok{"data"}\NormalTok{, }\StringTok{"gfdl{-}esm4{-}tempmax{-}ithaca.csv"}\NormalTok{), DataFrame, delim}\OperatorTok{=}\StringTok{","}\NormalTok{)}
\end{Highlighting}
\end{Shaded}

\begin{tabular}{r|cc}
    & Day & TempMax\\
    \hline
    & Date & Float64\\
    \hline
    1 & 1850-01-01 & 2.33492 \\
    2 & 1850-01-02 & 5.37603 \\
    3 & 1850-01-03 & 2.89517 \\
    4 & 1850-01-04 & -5.93915 \\
    5 & 1850-01-05 & -3.71647 \\
    6 & 1850-01-06 & -5.10975 \\
    7 & 1850-01-07 & -2.69068 \\
    8 & 1850-01-08 & 0.0958496 \\
    9 & 1850-01-09 & 0.242303 \\
    10 & 1850-01-10 & -0.239172 \\
    11 & 1850-01-11 & -1.82245 \\
    12 & 1850-01-12 & -4.21055 \\
    13 & 1850-01-13 & -3.57349 \\
    14 & 1850-01-14 & -1.52216 \\
    15 & 1850-01-15 & -8.84281 \\
    16 & 1850-01-16 & 0.163324 \\
    17 & 1850-01-17 & 0.697412 \\
    18 & 1850-01-18 & 0.0866333 \\
    19 & 1850-01-19 & -1.69102 \\
    20 & 1850-01-20 & -8.24766 \\
    21 & 1850-01-21 & -2.60285 \\
    22 & 1850-01-22 & 2.38784 \\
    23 & 1850-01-23 & 4.52969 \\
    24 & 1850-01-24 & 9.13345 \\
    $\dots$ & $\dots$ & $\dots$ \\
\end{tabular}

First, we want to find the annual maxima.

\begin{Shaded}
\begin{Highlighting}[]
\NormalTok{dat.Year }\OperatorTok{=} \FunctionTok{year}\NormalTok{.(dat.Day)}
\NormalTok{dat\_gp }\OperatorTok{=} \PreprocessorTok{@groupby}\NormalTok{(dat, }\OperatorTok{:}\NormalTok{Year)}
\NormalTok{dat\_max }\OperatorTok{=} \PreprocessorTok{@combine}\NormalTok{(dat\_gp, }\OperatorTok{:}\NormalTok{AnnMax}\OperatorTok{=}\FunctionTok{maximum}\NormalTok{(}\OperatorTok{:}\NormalTok{TempMax))}
\FunctionTok{plot}\NormalTok{(dat\_max.Year, dat\_max.AnnMax, lw}\OperatorTok{=}\FloatTok{2}\NormalTok{, marker}\OperatorTok{=:}\NormalTok{o, markersize}\OperatorTok{=}\FloatTok{3}\NormalTok{, xlabel}\OperatorTok{=}\StringTok{"Year"}\NormalTok{, ylabel}\OperatorTok{=}\StringTok{"Annual Maximum Temperature (°C)"}\NormalTok{)}
\end{Highlighting}
\end{Shaded}

\pandocbounded{\includegraphics[keepaspectratio]{hw03_files/mediabag/hw03_files/figure-pdf/cell-16-output-1.pdf}}

\paragraph{Problem 3.2}\label{problem-3.2}

Now, let's fit a stationary GEV model.

\begin{Shaded}
\begin{Highlighting}[]
\KeywordTok{function} \FunctionTok{gev\_stat\_loglik}\NormalTok{(p, x) }
\NormalTok{    u, s, xi }\OperatorTok{=}\NormalTok{ p}
\NormalTok{    ll }\OperatorTok{=} \FunctionTok{sum}\NormalTok{(}\FunctionTok{logpdf}\NormalTok{.(}\FunctionTok{GeneralizedExtremeValue}\NormalTok{(u, s, xi), x))}
    \ControlFlowTok{return}\NormalTok{ ll}
\KeywordTok{end}

\NormalTok{lb }\OperatorTok{=}\NormalTok{ [}\FloatTok{0.0}\NormalTok{, }\FloatTok{0.1}\NormalTok{, }\OperatorTok{{-}}\FloatTok{1.0}\NormalTok{]}
\NormalTok{ub }\OperatorTok{=}\NormalTok{ [}\FloatTok{50.0}\NormalTok{, }\FloatTok{10.0}\NormalTok{, }\FloatTok{1.0}\NormalTok{]}
\NormalTok{x₀ }\OperatorTok{=}\NormalTok{ [}\FloatTok{30.0}\NormalTok{, }\FloatTok{5.0}\NormalTok{, }\FloatTok{0.0}\NormalTok{]}
\NormalTok{optim\_out }\OperatorTok{=}\NormalTok{ Optim.}\FunctionTok{optimize}\NormalTok{(p }\OperatorTok{{-}\textgreater{}} \FunctionTok{{-}gev\_stat\_loglik}\NormalTok{(p, dat\_max.AnnMax), lb, ub, x₀)}
\NormalTok{p\_stat }\OperatorTok{=}\NormalTok{ optim\_out.minimizer}
\PreprocessorTok{@show} \FunctionTok{round}\NormalTok{.(p\_stat; digits}\OperatorTok{=}\FloatTok{1}\NormalTok{);}
\end{Highlighting}
\end{Shaded}

\begin{verbatim}
round.(p_stat; digits = 1) = [27.8, 1.0, -0.2]
\end{verbatim}

This is a Weibull distribution due to the negative shape parameter
\(\xi\), and the 100-year return level is the 0.99 quantile, or
30.8\(^\circ C\).

Let's plot the fitted distribution and the data.

\begin{Shaded}
\begin{Highlighting}[]
\FunctionTok{histogram}\NormalTok{(dat\_max.AnnMax, normalize}\OperatorTok{=:}\NormalTok{pdf, label}\OperatorTok{=}\StringTok{"Data"}\NormalTok{, xlabel}\OperatorTok{=}\StringTok{"Annual Maximal Temperature (°C)"}\NormalTok{, ylabel}\OperatorTok{=}\StringTok{"PDF"}\NormalTok{, color}\OperatorTok{=:}\NormalTok{lightblue)}
\FunctionTok{plot!}\NormalTok{(}\FunctionTok{GeneralizedExtremeValue}\NormalTok{(p\_stat[}\FloatTok{1}\NormalTok{], p\_stat[}\FloatTok{2}\NormalTok{], p\_stat[}\FloatTok{3}\NormalTok{]), color}\OperatorTok{=:}\NormalTok{red, label}\OperatorTok{=}\StringTok{"Fitted Stationary GEV"}\NormalTok{, lw}\OperatorTok{=}\FloatTok{2}\NormalTok{)}
\end{Highlighting}
\end{Shaded}

\begin{figure}[H]

\centering{

\pandocbounded{\includegraphics[keepaspectratio]{hw03_files/mediabag/hw03_files/figure-pdf/fig-stat-gev-output-1.pdf}}

}

\caption{\label{fig-stat-gev}}

\end{figure}%

\paragraph{Problem 3.3}\label{problem-3.3}

Now, let's fit a nonstationary GEV, where the location parameter \(\mu\)
is a function of time, \emph{e.g.} \(\mu = \mu_0 + \mu_1 t\).

\begin{Shaded}
\begin{Highlighting}[]
\KeywordTok{function} \FunctionTok{gev\_nonstat\_loglik}\NormalTok{(p, x, yr) }
\NormalTok{    u₀, u₁, s, xi }\OperatorTok{=}\NormalTok{ p}
\NormalTok{    t }\OperatorTok{=}\NormalTok{ yr }\OperatorTok{.{-}}\NormalTok{ yr[}\FloatTok{1}\NormalTok{]}
\NormalTok{    u }\OperatorTok{=}\NormalTok{ u₀ }\OperatorTok{.+}\NormalTok{ u₁ }\OperatorTok{.*}\NormalTok{ t}
\NormalTok{    ll }\OperatorTok{=} \FunctionTok{sum}\NormalTok{(}\FunctionTok{logpdf}\NormalTok{.(}\FunctionTok{GeneralizedExtremeValue}\NormalTok{.(u, s, xi), x))}
    \ControlFlowTok{return}\NormalTok{ ll}
\KeywordTok{end}

\NormalTok{lb }\OperatorTok{=}\NormalTok{ [}\FloatTok{0.0}\NormalTok{, }\FloatTok{0.0}\NormalTok{, }\FloatTok{0.1}\NormalTok{, }\OperatorTok{{-}}\FloatTok{1.0}\NormalTok{]}
\NormalTok{ub }\OperatorTok{=}\NormalTok{ [}\FloatTok{40.0}\NormalTok{, }\FloatTok{1.0}\NormalTok{, }\FloatTok{10.0}\NormalTok{, }\FloatTok{1.0}\NormalTok{]}
\NormalTok{x₀ }\OperatorTok{=}\NormalTok{ [}\FloatTok{10.0}\NormalTok{, }\FloatTok{0.5}\NormalTok{, }\FloatTok{5.0}\NormalTok{, }\FloatTok{0.0}\NormalTok{]}
\NormalTok{optim\_out }\OperatorTok{=}\NormalTok{ Optim.}\FunctionTok{optimize}\NormalTok{(p }\OperatorTok{{-}\textgreater{}} \FunctionTok{{-}gev\_nonstat\_loglik}\NormalTok{(p, dat\_max.AnnMax, dat\_max.Year), lb, ub, x₀)}
\NormalTok{p\_nonstat }\OperatorTok{=}\NormalTok{ optim\_out.minimizer}
\PreprocessorTok{@show} \FunctionTok{round}\NormalTok{.(p\_nonstat; digits}\OperatorTok{=}\FloatTok{1}\NormalTok{);}
\end{Highlighting}
\end{Shaded}

\begin{verbatim}
round.(p_nonstat; digits = 1) = [27.7, 0.0, 1.0, -0.2]
\end{verbatim}

We want to calculate and plot how the 100-year return level changed over
time according to this model. We also want to compare to the 100-year
return level from Problem 3.3

\begin{Shaded}
\begin{Highlighting}[]
\NormalTok{u }\OperatorTok{=} \FunctionTok{collect}\NormalTok{(p\_nonstat[}\FloatTok{1}\NormalTok{] }\OperatorTok{.+}\NormalTok{ p\_nonstat[}\FloatTok{2}\NormalTok{] }\OperatorTok{*}\NormalTok{ (}\FloatTok{1}\OperatorTok{:}\FunctionTok{nrow}\NormalTok{(dat\_max)))}
\NormalTok{nonstat\_rl }\OperatorTok{=} \FunctionTok{quantile}\NormalTok{.(}\FunctionTok{GeneralizedExtremeValue}\NormalTok{.(u, p\_nonstat[}\FloatTok{3}\NormalTok{], p\_nonstat[}\FloatTok{4}\NormalTok{]), }\FloatTok{0.99}\NormalTok{)}

\FunctionTok{plot}\NormalTok{(dat\_max.Year, nonstat\_rl, color}\OperatorTok{=:}\NormalTok{black, lw}\OperatorTok{=}\FloatTok{2}\NormalTok{, xlabel}\OperatorTok{=}\StringTok{"Year"}\NormalTok{, ylabel}\OperatorTok{=}\StringTok{"100{-}Year Return Level"}\NormalTok{, label}\OperatorTok{=}\StringTok{"Nonstationary Model"}\NormalTok{)}
\FunctionTok{hline!}\NormalTok{([}\FunctionTok{quantile}\NormalTok{(}\FunctionTok{GeneralizedExtremeValue}\NormalTok{(p\_stat[}\FloatTok{1}\NormalTok{], p\_stat[}\FloatTok{2}\NormalTok{], p\_stat[}\FloatTok{3}\NormalTok{]), }\FloatTok{0.99}\NormalTok{)], color}\OperatorTok{=:}\NormalTok{red, linestyle}\OperatorTok{=:}\NormalTok{dash, lw}\OperatorTok{=}\FloatTok{2}\NormalTok{, label}\OperatorTok{=}\StringTok{"Stationary Model"}\NormalTok{)}
\end{Highlighting}
\end{Shaded}

\pandocbounded{\includegraphics[keepaspectratio]{hw03_files/mediabag/hw03_files/figure-pdf/cell-20-output-1.pdf}}

The non-stationary return levels increase from about 30.6\(^\circ C\) to
closer to 31\(^\circ C\), while the stationary return level is about the
average of that range.

\subsubsection{Problem 4}\label{problem-4}

\paragraph{Problem 4.1}\label{problem-4.1}

We don't need to re-load the dataset, so let's find the exceedances over
28\(^\circ\) C.

\begin{Shaded}
\begin{Highlighting}[]
\NormalTok{temp\_high }\OperatorTok{=}\NormalTok{ dat[dat.TempMax }\OperatorTok{.\textgreater{}} \FloatTok{28}\NormalTok{, }\OperatorTok{:}\NormalTok{]}
\end{Highlighting}
\end{Shaded}

\begin{tabular}{r|ccc}
    & Day & TempMax & Year\\
    \hline
    & Date & Float64 & Int64\\
    \hline
    1 & 1853-08-13 & 28.734 & 1853 \\
    2 & 1854-07-12 & 29.5106 & 1854 \\
    3 & 1854-07-18 & 28.131 & 1854 \\
    4 & 1854-07-19 & 28.9641 & 1854 \\
    5 & 1854-07-31 & 28.1494 & 1854 \\
    6 & 1854-08-01 & 28.2747 & 1854 \\
    7 & 1854-08-02 & 28.6134 & 1854 \\
    8 & 1854-08-03 & 28.7721 & 1854 \\
    9 & 1857-07-07 & 28.0677 & 1857 \\
    10 & 1858-07-29 & 28.0767 & 1858 \\
    11 & 1860-07-01 & 28.5655 & 1860 \\
    12 & 1860-07-29 & 29.743 & 1860 \\
    13 & 1860-08-08 & 28.4803 & 1860 \\
    14 & 1861-06-30 & 28.3882 & 1861 \\
    15 & 1861-07-20 & 29.0316 & 1861 \\
    16 & 1861-07-21 & 28.5016 & 1861 \\
    17 & 1861-07-22 & 29.5198 & 1861 \\
    18 & 1861-08-10 & 28.1314 & 1861 \\
    19 & 1861-08-17 & 29.2906 & 1861 \\
    20 & 1862-07-16 & 29.7026 & 1862 \\
    21 & 1863-07-11 & 28.8037 & 1863 \\
    22 & 1864-06-18 & 28.2663 & 1864 \\
    23 & 1864-07-15 & 28.0548 & 1864 \\
    24 & 1864-08-10 & 28.0876 & 1864 \\
    $\dots$ & $\dots$ & $\dots$ & $\dots$ \\
\end{tabular}

To decluster, we calculate the extremal index using the estimator
\[\hat{\theta}(u) = \frac{2\left(\sum_{i-1}^{N-1} T_i\right)^2}{(N-1)\sum_{i=1}^{N-1}T_i^2}.\]

\begin{Shaded}
\begin{Highlighting}[]
\NormalTok{S }\OperatorTok{=} \FunctionTok{findall}\NormalTok{(dat.TempMax }\OperatorTok{.\textgreater{}} \FloatTok{28}\NormalTok{)}
\NormalTok{N }\OperatorTok{=} \FunctionTok{length}\NormalTok{(S)}
\NormalTok{T }\OperatorTok{=} \FunctionTok{diff}\NormalTok{(S)}
\NormalTok{θ }\OperatorTok{=} \FloatTok{2} \OperatorTok{*} \FunctionTok{sum}\NormalTok{(T)}\OperatorTok{\^{}}\FloatTok{2} \OperatorTok{/}\NormalTok{ ((N}\OperatorTok{{-}}\FloatTok{1}\NormalTok{) }\OperatorTok{*} \FunctionTok{sum}\NormalTok{(T}\OperatorTok{.\^{}}\FloatTok{2}\NormalTok{))}
\end{Highlighting}
\end{Shaded}

\begin{verbatim}
0.47832378879687915
\end{verbatim}

This means that the declustering time is 2.0 days. Now we want to assign
data that is within that period to appropriate clusters. Let's write a
function to do that and then take the maximum value of the exceedances
within each cluster.

\begin{Shaded}
\begin{Highlighting}[]
\CommentTok{\# cluster data points which occur within period}
\KeywordTok{function} \FunctionTok{assign\_cluster}\NormalTok{(dat, period)}
\NormalTok{    cluster\_index }\OperatorTok{=} \FloatTok{1}
\NormalTok{    clusters }\OperatorTok{=} \FunctionTok{zeros}\NormalTok{(}\DataTypeTok{Int}\NormalTok{, }\FunctionTok{length}\NormalTok{(dat))}
    \ControlFlowTok{for}\NormalTok{ i }\KeywordTok{in} \FloatTok{1}\OperatorTok{:}\FunctionTok{length}\NormalTok{(dat)}
        \ControlFlowTok{if}\NormalTok{ clusters[i] }\OperatorTok{==} \FloatTok{0}
\NormalTok{            clusters[}\FunctionTok{findall}\NormalTok{(}\FunctionTok{abs}\NormalTok{.(dat }\OperatorTok{.{-}}\NormalTok{ dat[i]) }\OperatorTok{.\textless{}=}\NormalTok{ period)] }\OperatorTok{.=}\NormalTok{ cluster\_index}
\NormalTok{            cluster\_index }\OperatorTok{+=} \FloatTok{1}
        \ControlFlowTok{end}
    \ControlFlowTok{end}
    \ControlFlowTok{return}\NormalTok{ clusters}
\KeywordTok{end}

\CommentTok{\# cluster exceedances that occur within a two{-}day window}
\CommentTok{\# @transform is a macro from DataFramesMeta.jl which adds a new column based on a data transformation}
\NormalTok{temp\_high }\OperatorTok{=} \PreprocessorTok{@transform}\NormalTok{ temp\_high }\OperatorTok{:}\NormalTok{cluster }\OperatorTok{=} \FunctionTok{assign\_cluster}\NormalTok{(}\OperatorTok{:}\NormalTok{Day, }\BuiltInTok{Dates}\NormalTok{.}\FunctionTok{Day}\NormalTok{(}\FloatTok{2}\NormalTok{))}
\CommentTok{\# find maximum value within cluster}
\NormalTok{temp\_decluster }\OperatorTok{=} \FunctionTok{combine}\NormalTok{(temp\_high }\OperatorTok{{-}\textgreater{}}\NormalTok{ temp\_high[}\FunctionTok{argmax}\NormalTok{(temp\_high.TempMax), }\OperatorTok{:}\NormalTok{], }\FunctionTok{groupby}\NormalTok{(temp\_high, }\OperatorTok{:}\NormalTok{cluster))}
\end{Highlighting}
\end{Shaded}

\begin{tabular}{r|cccc}
    & cluster & Day & TempMax & Year\\
    \hline
    & Int64 & Date & Float64 & Int64\\
    \hline
    1 & 1 & 1853-08-13 & 28.734 & 1853 \\
    2 & 2 & 1854-07-12 & 29.5106 & 1854 \\
    3 & 3 & 1854-07-19 & 28.9641 & 1854 \\
    4 & 4 & 1854-07-31 & 28.1494 & 1854 \\
    5 & 5 & 1854-08-03 & 28.7721 & 1854 \\
    6 & 6 & 1857-07-07 & 28.0677 & 1857 \\
    7 & 7 & 1858-07-29 & 28.0767 & 1858 \\
    8 & 8 & 1860-07-01 & 28.5655 & 1860 \\
    9 & 9 & 1860-07-29 & 29.743 & 1860 \\
    10 & 10 & 1860-08-08 & 28.4803 & 1860 \\
    11 & 11 & 1861-06-30 & 28.3882 & 1861 \\
    12 & 12 & 1861-07-22 & 29.5198 & 1861 \\
    13 & 13 & 1861-08-10 & 28.1314 & 1861 \\
    14 & 14 & 1861-08-17 & 29.2906 & 1861 \\
    15 & 15 & 1862-07-16 & 29.7026 & 1862 \\
    16 & 16 & 1863-07-11 & 28.8037 & 1863 \\
    17 & 17 & 1864-06-18 & 28.2663 & 1864 \\
    18 & 18 & 1864-07-15 & 28.0548 & 1864 \\
    19 & 19 & 1864-08-11 & 30.1618 & 1864 \\
    20 & 20 & 1864-08-31 & 28.8849 & 1864 \\
    21 & 21 & 1865-08-18 & 28.0744 & 1865 \\
    22 & 22 & 1867-08-11 & 30.3048 & 1867 \\
    23 & 23 & 1868-07-21 & 28.1273 & 1868 \\
    24 & 24 & 1868-07-24 & 28.0214 & 1868 \\
    $\dots$ & $\dots$ & $\dots$ & $\dots$ & $\dots$ \\
\end{tabular}

Now we want to count the number of exceedances by year and plot how
these occurrences are changing over time.

\begin{Shaded}
\begin{Highlighting}[]
\NormalTok{temp\_decluster.Year }\OperatorTok{=} \FunctionTok{year}\NormalTok{.(temp\_decluster.Day)}
\NormalTok{temp\_count }\OperatorTok{=} \FunctionTok{combine}\NormalTok{(}\FunctionTok{groupby}\NormalTok{(temp\_decluster, }\OperatorTok{:}\NormalTok{Year), nrow)}
\FunctionTok{bar}\NormalTok{(temp\_count.Year, temp\_count.nrow, label}\OperatorTok{=}\StringTok{"Number of Clusters"}\NormalTok{)}
\end{Highlighting}
\end{Shaded}

\pandocbounded{\includegraphics[keepaspectratio]{hw03_files/mediabag/hw03_files/figure-pdf/cell-24-output-1.pdf}}

Visually, it appears as though there has been an uptick in the number of
clusters post-2000, with more years containing 3+ clusters. This would
be hard to pull out without a model that depends on a broader warming
signal, since the most of the data record looks stationary, so it's
unclear whether this is actually a trend. As a result, we might use
stationary Poisson process for the occurrences, or we could try to fit
one that was based on a covariate like global mean temperature.

\paragraph{Problem 4.2}\label{problem-4.2}

Let's now fit a stationary GPD to see what how the distribution of
conditional exceedances looks.

\begin{Shaded}
\begin{Highlighting}[]
\NormalTok{temp\_exceedances }\OperatorTok{=}\NormalTok{ temp\_high.TempMax }\OperatorTok{.{-}} \FloatTok{28}
\NormalTok{gpd\_fit }\OperatorTok{=}\NormalTok{ Optim.}\FunctionTok{optimize}\NormalTok{(p }\OperatorTok{{-}\textgreater{}} \FunctionTok{{-}sum}\NormalTok{(}\FunctionTok{logpdf}\NormalTok{(}\FunctionTok{GeneralizedPareto}\NormalTok{(}\FloatTok{0.0}\NormalTok{, p[}\FloatTok{1}\NormalTok{], p[}\FloatTok{2}\NormalTok{]), temp\_exceedances)), [}\FloatTok{0.0}\NormalTok{, }\OperatorTok{{-}}\ConstantTok{Inf}\NormalTok{], [}\ConstantTok{Inf}\NormalTok{, }\ConstantTok{Inf}\NormalTok{], [}\FloatTok{1.0}\NormalTok{, }\FloatTok{1.0}\NormalTok{]).minimizer}
\end{Highlighting}
\end{Shaded}

\begin{verbatim}
2-element Vector{Float64}:
  0.8510410229847296
 -0.18775785233967718
\end{verbatim}

Let's visualize this distribution (Figure~\ref{fig-gpd-fit}).

\begin{Shaded}
\begin{Highlighting}[]
\FunctionTok{histogram}\NormalTok{(temp\_decluster.TempMax, normalize}\OperatorTok{=:}\NormalTok{pdf, label}\OperatorTok{=}\StringTok{"Data"}\NormalTok{, xlabel}\OperatorTok{=}\StringTok{"Conditional Excess Distribution (°C)"}\NormalTok{, ylabel}\OperatorTok{=}\StringTok{"PDF"}\NormalTok{, color}\OperatorTok{=:}\NormalTok{lightblue)}
\FunctionTok{plot!}\NormalTok{(}\FunctionTok{GeneralizedPareto}\NormalTok{(}\FloatTok{0.0}\NormalTok{, gpd\_fit[}\FloatTok{1}\NormalTok{], gpd\_fit[}\FloatTok{2}\NormalTok{]) }\OperatorTok{.+} \FloatTok{28}\NormalTok{, color}\OperatorTok{=:}\NormalTok{red, label}\OperatorTok{=}\StringTok{"Fitted Stationary GPD"}\NormalTok{, lw}\OperatorTok{=}\FloatTok{2}\NormalTok{)}
\end{Highlighting}
\end{Shaded}

\begin{figure}[H]

\centering{

\pandocbounded{\includegraphics[keepaspectratio]{hw03_files/mediabag/hw03_files/figure-pdf/fig-gpd-fit-output-1.pdf}}

}

\caption{\label{fig-gpd-fit}GPD fit for temperature exceedances over
28°C.}

\end{figure}%

Now, let's fit a Poisson model for the number of exceedances in each
year.

\begin{Shaded}
\begin{Highlighting}[]
\NormalTok{poisson\_fit }\OperatorTok{=}\NormalTok{ Optim.}\FunctionTok{optimize}\NormalTok{(p }\OperatorTok{{-}\textgreater{}} \FunctionTok{{-}sum}\NormalTok{(}\FunctionTok{logpdf}\NormalTok{.(}\FunctionTok{Poisson}\NormalTok{(p[}\FloatTok{1}\NormalTok{]), temp\_count.nrow)), [}\FloatTok{0.1}\NormalTok{], [}\FloatTok{10.0}\NormalTok{], [}\FloatTok{2.0}\NormalTok{]).minimizer}
\PreprocessorTok{@show} \FunctionTok{round}\NormalTok{(poisson\_fit[}\FloatTok{1}\NormalTok{]; digits}\OperatorTok{=}\FloatTok{2}\NormalTok{);}
\end{Highlighting}
\end{Shaded}

\begin{verbatim}
round(poisson_fit[1]; digits = 2) = 1.73
\end{verbatim}

The 100-year return level is given by the solution to the quantile
equation,

\[
\begin{aligned}
\text{RL}_{100} &= u + \frac{\sigma}{\xi}\left((100\lambda)^\xi - 1\right)
\\ &= 28 - \frac{0.85}{0.19}\left(100 * 1.73)^{-0.19} - 1\right) \\
&= 30.8^\circ \text{C}
\end{aligned}
\]

This is effectively the same as the result from the stationary GEV
analysis in Problem 3 (the difference is down to the second decimal
place).

\subsection{References}\label{references}




\end{document}
